\documentclass{article}
\usepackage[T1]{fontenc}
\usepackage{lmodern}
\usepackage{graphicx}
\usepackage{amsmath,amssymb}
\usepackage[a4paper,margin=2cm]{geometry}
\usepackage[absolute,overlay]{textpos}
\setlength{\TPHorizModule}{1mm}
\setlength{\TPVertModule}{1mm}
\usepackage[indonesian]{babel}
\usepackage{hyperref}
\setlength{\emergencystretch}{2em}
% unit: Halaman 38

\begin{document}

% === Halaman 38 (llm) ===

\section*{Penggunaan Kurva Eliptik di dalam Kriptografi}

\begin{itemize}
    \item Bagian inti dari sistem kriptografi kunci-publik yang melibatkan kurva eliptik adalah \textbf{grup eliptik} (himpunan titik-titik pada kurva eliptik dan sebuah operasi biner $+$).
    \item Operasi matematika yang mendasari:
    \begin{itemize}
        \item Jika RSA mempunyai operasi perpangkatan sebagai operasi matematika yang mendasarinya, maka
        \item ECC memiliki operasi perkalian titik (kP)
    \end{itemize}
\end{itemize}

\vspace{20mm}

\begin{flushright}
\small{*) Sumber bahan: \textcolor{red}{Debdeep Mukhopadhyay, \textcolor{red}{Elliptic Curve Cryptography}, \textcolor{red}{Dept of Computer Sc and Engg IIT Madras}}
\end{flushright}

\end{document}
